\section{Finite matrix certificates and an expanding example}
\label{sec:mc-section}

The component phases in \cref{thm:dc-complex} have a finite algebraic
counterpart.  Its witnesses may have growing dimension while the code
being bounded retains any chosen block width.  This distinction permits
a finite-field result without a spectral limit of recovery matrices.
Central layers then transfer its bound to a specified class of nonlinear
finite-group recovery rules.

\begin{lemma}[Rank identity]
\formal{RankCertificate.lean}{rank_identity}
\label{lem:mc-identity}
For endomorphisms $A,L,R$ of any finite-dimensional vector space over any
field $F$,
\begin{equation}
 2\operatorname{rank}(I-A)
 \ge\operatorname{rank}(L-R)-\operatorname{rank}(AL-RA).
 \label{eq:mc-rank-identity}
\end{equation}
No invertibility or commutativity is required.
\end{lemma}
\begin{outline}
For $M=I-A$ the exact identity is
$L-R=(AL-RA)+ML-RM$.  Apply rank subadditivity and
$\operatorname{rank}(ML),\operatorname{rank}(RM)\le\operatorname{rank}M$.
\end{outline}

\begin{proposition}[Matrix-label certificate]
\formal{MatrixCertificate.lean}{matrix_label_certificate,matrix_label_certificate_extension}
\label{prop:mc-certificate}
Let $A$ have $t\times t$ blocks indexed by $N$ coordinates, supported on
a paired incidence relation $E$.  Give its ports arbitrary labels
$X_p,Y_q\in\operatorname{Mat}_m(F)$, where $m\ge1$.  Then
\begin{equation}
 \operatorname{rank}_F(I-A)\ge
 \frac{t\sum_{p}\operatorname{rank}_F(X_p-Y_p)
 -\sum_{L_pR_q\in E}\operatorname{rank}_F(A_{qp})
                         \operatorname{rank}_F(X_p-Y_q)}{2m}.
 \label{eq:mc-certificate-equation}
\end{equation}
Labels over a field extension of $F$ are also valid witnesses, with
the ranks of the original matrices unchanged.
\end{proposition}
\begin{outline}
Use $\widehat A=A\otimes I_m$ and block-diagonal operators
$L_{pp}=I_t\otimes X_p$, $R_{qq}=I_t\otimes Y_q$.
Then
\[
 \operatorname{rank}(I-\widehat A)=m\operatorname{rank}(I-A),\qquad
 \operatorname{rank}(L-R)=t\sum_p\operatorname{rank}(X_p-Y_p),
\]
and the $(q,p)$ block of $\widehat A L-R\widehat A$ is
$A_{qp}\otimes(X_p-Y_q)$.  Tensor-product ranks multiply over any field.
Writing a block matrix as the sum of its single-block inclusions bounds
its rank by the sum of their ranks.  \cref{lem:mc-identity} gives the
displayed inequality. Scalar extension preserves matrix rank, so the same
argument applies to labels over any extension field.
\end{outline}

In particular, full-rank differences on matched ports and low-rank
differences on incidence edges give a useful certificate.  The following
construction shows that this situation is compatible with uniform
expansion and large matched distances.

\begin{theorem}[An existential characteristic-two family]
\notformalized
\label{thm:mc-expanding-family}
There are constants $\Delta\ge1$, $h>0$, and $m_0$ such that for every
integer $m\ge m_0$ there is a simple paired incidence graph $E_m$ on
$2N_m$ ports, where
\[
 N_m=2^{m^2}(2^m-1)^2,
\]
with the following properties:
\begin{enumerate}
\item Its maximum degree is at most $\Delta$, and every port set $S$ with
      $|S|\le N_m$ has at least $h|S|$ edges leaving it.
\item Every matched pair has distance at least $m$.
\item For every field $F$ of characteristic two, every integer $t\ge1$,
      and every recovery matrix $A$ of block width $t$ supported on $E_m$,
\begin{equation}
 \operatorname{rank}_F(I-A)
 \ge\frac{tN_m}{2}
       -\frac1{2m}\sum_{L_pR_q\in E_m}\operatorname{rank}_F(A_{qp})
 \ge\left(\frac12-\frac{\Delta}{2m}\right)tN_m.
 \label{eq:mc-expander-rank}
\end{equation}
\end{enumerate}
The graph is chosen independently of $F,t,A$.  In particular the bound
is uniform in $t$, even if $t$ grows with $m$.
\end{theorem}

\begin{informalproof}
\emph{An outer graph with rank-one steps.}
Put $s=2^m$ and $D_m=(s-1)^2$.  On the additive group
$\operatorname{Mat}_m(\mathbb F_2)$ join $B$ to $B+uv^{\mathsf T}$
for each ordered pair of nonzero vectors $u,v\in\mathbb F_2^m$.
Each nonzero rank-one matrix has exactly one such representation over
$\mathbb F_2$, so this is a simple $D_m$-regular graph $G_m$.
The edge label $(u,v)$ is unchanged by reversal.

For a matrix $C$ of rank $r$, the additive character
$B\mapsto(-1)^{\sum_{i,j}C_{ij}B_{ij}}$ has normalized eigenvalue
\begin{equation}
 \lambda_r=
 \frac{2^{2m-r}-2^{m+1}+1}{(2^m-1)^2}.
 \label{eq:mc-outer-spectrum}
\end{equation}
Indeed, among the $s-1$ nonzero vectors $u$, exactly $2^{m-r}-1$
satisfy $u^{\mathsf T}C=0$.  Summing the character over nonzero $v$
gives $s-1$ for those $u$ and $-1$ for the others.  The resulting sum
is $2^{2m-r}-2s+1$.  These characters form an orthogonal eigenbasis.
For $m\ge2$ and $r\ge1$, the eigenvalues decrease from at most $1/2$
to $-1/(s-1)\ge-1/3$.  Thus all nonconstant eigenvalues have absolute
value at most $1/2$.

\emph{Reducing the degree.}
Fix $k=1024$.  Theorem~1 of Bordenave~\cite{Bordenave} implies that,
for every sufficiently large $D$, there is a simple $k$-regular graph
$H_D$ on $D$ vertices whose normalized nonconstant eigenvalues have
absolute value at most $1/10$.  In fact
$(2\sqrt{1023}+1)/1024<1/10$, and $kD$ is always even.
Choose one such graph on the label set of $G_m$.

Use the zig-zag construction of Reingold, Vadhan, and
Wigderson~\cite[Section~3]{RVW}.  Its vertex set consists of pairs $(B,a)$,
where $a$ is an outer edge label.  A step first changes $a$ along one
edge of $H_{D_m}$, then takes that labelled edge of $G_m$, and finally
changes the new label along one edge of $H_{D_m}$.
The resulting undirected multigraph $Z_m$ has $N_m$ vertices and degree
$\Delta=k^2$.  Every edge changes $B$ by a matrix of rank exactly one.
There are no loops; we permit parallel edges in the spectral accounting.

For completeness, the following loose spectral estimate suffices.
Let $P$ average within each label cloud, let $\mathcal H$ apply the
normalized adjacency of $H_{D_m}$ within clouds, and let $\mathcal R$
be the outer rotation, an orthogonal involution.
The normalized adjacency of $Z_m$ is
$\mathcal H\mathcal R\mathcal H$.
For a vector $f$ orthogonal to the constants, write
$f=f_\parallel+f_\perp$, where $f_\parallel=Pf$, and put
$e=\mathcal Hf_\perp$.  Then $\|e\|\le\|f_\perp\|/10$.
On cloud-constant vectors $P\mathcal RP$ is the outer normalized
adjacency.  Hence
\[
 \bigl|\langle f,\mathcal H\mathcal R\mathcal Hf\rangle\bigr|
 \le\tfrac12\|f_\parallel\|^2
      +\tfrac2{10}\|f_\parallel\|\|f_\perp\|
      +\tfrac1{100}\|f_\perp\|^2
 \le0.71\|f\|^2.
\]
Since the adjacency is self-adjoint, all its nonconstant eigenvalues
have absolute value at most $\lambda=0.71<1$.

\emph{Ports and their prescribed pairing.}
For a vertex $x=(B,a)$ put $\theta(x)=(B+I_m,a)$.
Take the bipartite double cover of $Z_m$, keep the physical left port
at $x$ named $L_x$, and name the physical right port at $\theta(y)$ by
$R_y$.  Finally remove parallel copies of edges to obtain the simple
graph $E_m$.  Thus
\begin{equation}
 L_xR_y\in E_m\quad\Longleftrightarrow\quad
 x\text{ and }\theta(y)\text{ are adjacent in }Z_m.
 \label{eq:mc-shifted-pairing}
\end{equation}
Give these ports the matrix labels $X_x=B_x$ and $Y_y=B_y+I_m$.
Each edge difference has rank one.  Each matched difference is $I_m$
and has rank $m$.  Along any port path the matrix differences telescope;
rank subadditivity therefore forces every matched path to have at least
$m$ edges.  In particular there are no self-demand edges when $m>1$.

The double cover has eigenvalues equal to the positive and negative
eigenvalues of $Z_m$.  Its eigenvalue $1$ is simple, and every other
eigenvalue is at most $\lambda$; the additional eigenvalue $-1$ is harmless
for this upper bound.  For a port set $S$ of size $q\le N_m$, apply
the spectral gap to $1_S-q/(2N_m)$.  It gives at least
$\Delta(1-\lambda)q/2$ crossing edges counted with multiplicity.
Every simple edge has multiplicity at most $\Delta$.  Deduplication
therefore leaves at least $hq$ crossing edges for
$h=(1-\lambda)/2>0$, and degree at most $\Delta$.

\emph{The rank bound.}
Interpret the binary labels $X_x,Y_y$ in any field of characteristic two.
Their ranks are preserved by extension of the prime field.
\cref{prop:mc-certificate} gives the first inequality in
\eqref{eq:mc-expander-rank}.  The second follows from
$\operatorname{rank}A_{yx}\le t$ and $|E_m|\le\Delta N_m$.
\end{informalproof}

\begin{corollary}[Linear codes on the expanding family]
\partformal{LinearRecovery.lean}{linear_recovery_matrix_exists,linear_code_rate_of_matrix_labels}{Construction of the expanding family and its full-rank matched and rank-one edge labels.}
\label{cor:mc-linear-codes}
Over any characteristic-two field, every linear storage space on $E_m$
with $t$ symbols per coordinate has normalized dimension at most
$1/2+\Delta/(2m)$.  This includes linear spaces initially supplied with
nonlinear recovery functions.
\end{corollary}
\begin{informalproof}
On a linear space, determination of a coordinate by its neighbours means
that the kernel of the neighbour projection is contained in the kernel
of the coordinate projection.  The latter therefore factors linearly
through the image of the former, and this linear map extends to the
ambient neighbour space.  Choose these linear recovery maps for all
coordinates to obtain $A$ with $\mathcal C\subseteq\ker(I-A)$, and use
\cref{thm:mc-expanding-family}.
\end{informalproof}

These graphs are not uniformly hyperfinite.  If deletion of $z$ ports
leaves components of size at most a fixed $K<N_m$, expansion summed
over the components gives
$h(2N_m-z)\le\Delta z$.  Thus every such separator has a fixed positive
linear size.  The preceding finite-field result consequently goes beyond
the hyperfinite hypothesis of \cref{cor:dc-hyperfinite}. The next theorem
also transfers it to nonlinear codes with group-expression recovery.

Two scope points are useful.  First, the binary translation $\theta$ is
an involutive automorphism of $Z_m$, so the demand relation in this
particular construction is symmetric.  Its contribution is a
characteristic-two bound uniform in block width, not an example requiring
asymmetric recovery.  Second, the inner graphs were chosen by an existence
theorem.  No efficient uniform generator or native tape realization is
asserted.  These are abstract relations; their degree, expansion, and
paired-port separation do not identify them with all native cones.

\subsection{From vector-linear bounds to group-expression recovery}

Fix a finite loopless recovery relation $E$ on $N>0$ coordinates and a
prime $p$. For a nontrivial finite $p$-group $G$, an \emph{endomorphic
recovery expression} uses permitted parent variables, constants in $G$,
multiplication, inversion, and fixed group endomorphisms $G\to G$.
Expressions are finite, with arbitrary nesting and repeated variables.
Each row has one fixed expression on the whole code. The code itself
can be any nonempty subset of $G^V$ satisfying these equations; it need
not be a subgroup. Arbitrary symbol maps inside an expression are not
included in this language.

\begin{theorem}[Central-layer transfer and the exact benchmark]
\notformalized
\label{thm:mc-group-transfer}
Suppose that for every integer $t\ge1$ and every $t\times t$ block
matrix $A$ over $\F_p$ supported on $E$,
$\dim\ker(I-A)\le\rho t$. Then every endomorphic-expression code
over every finite $p$-group satisfies $|\mathcal C|\le|G|^\rho$.
In particular, define
\[
 \gamma_p(E)=\sup_{t\ge1}\ max_{A\text{ supported on }E}
                    \frac{\dim_{\F_p}\ker(I-A)}{tN}.
\]
The supremum of $N^{-1}\log_{|G|}|\mathcal C|$ over all nontrivial
finite $p$-groups and their nonempty endomorphic-expression codes is
exactly $\gamma_p(E)$, the all-width vector-linear benchmark on $E$.
\end{theorem}
\begin{informalproof}
Use the lower exponent-$p$ central series
$P_1=G$, $P_{j+1}=P_j^p[P_j,G]$, where products denote generated
subgroups. Its terms are fully invariant. Each
$V_j=P_j/P_{j+1}$ is central in $G/P_{j+1}$ and is a vector space
over $\F_p$, of dimension $t_j$. Every endomorphism of $G$ therefore
induces a linear map on $V_j$.

Here is an elementary termination argument. Every nonzero finite
$\F_p$-representation of $G$ has a nonzero fixed vector: its additive
set has size divisible by $p$, every nontrivial orbit has size divisible
by $p$, and zero is fixed, so there are at least $p$ fixed vectors.
Iterate this observation on quotient representations to obtain a full
flag with trivial one-dimensional factors. On the regular module
$\F_p[G]$ of dimension $D=|G|$, every $g-1$ lowers this flag.
Thus its augmentation ideal $I$ satisfies $I^D=0$.
Inductively, $g\in P_j$ implies $g-1\in I^j$: the assertion is
preserved under products and inverses, powers satisfy
$g^p-1=(g-1)^p\in I^{j+1}$, and
\[
 [g,h]-1=g^{-1}h^{-1}
       \bigl((g-1)(h-1)-(h-1)(g-1)\bigr)\in I^{j+1}.
\]
Consequently $P_D=1$ and $\sum_jt_j=\log_p|G|$.

Project the actual code to $(G/P_j)^V$. Fix one of its attained
quotient assignments and choose lifts $r_v\in G/P_{j+1}$.
Every further lift is uniquely $r_va_v$, with $a_v\in V_j$.
Structural induction on an allowed expression $w$ gives
\[
 w((r_va_v)_v)=w((r_v)_v)L_w(a),
\]
where $L_w$ is linear over $\F_p$ and uses only the variables of $w$.
Variables give projections and constants give zero. Products add the
linear maps, inversion negates them, and an endomorphism composes with
its induced map on $V_j$. Centrality permits these operations even
when the base values do not commute. In this language $L_w$ is
independent of the chosen base assignment.

The recovery equations in the fibre therefore have the form
$(I-A_j)a=b$, with $t_j\times t_j$ blocks supported on $E$.
Their solutions form an affine translate of $\ker(I-A_j)$ when
nonempty. The actual code fibre can be a proper subset of this set;
no assertion that all quotient solutions lift is needed.
Multiply the uniform bounds on these successive fibres to get
\[
 \log_p|\mathcal C|\le\sum_j\dim\ker(I-A_j)
                      \le\rho\sum_jt_j=\rho\log_p|G|.
\]
Zero-width layers are omitted. Trivial groups and empty codes satisfy
the cardinality assertion separately.

Taking $\rho=N\gamma_p(E)$ proves the benchmark upper bound without
assuming its supremum is attained. Conversely take the additive group
$G=(\F_p)^t$. Every block linear map is an allowed endomorphism,
and sums of them realize each kernel in the definition of $\gamma_p$.
Thus every rate in that supremum is realized by the allowed language.
Finally any vector-linear storage space, even with initially nonlinear
decoders, admits linear decoders by factoring its parent projections,
as in \cref{cor:mc-linear-codes}. This identifies the benchmark with
the usual vector-linear one.
\end{informalproof}

Central lifting for ordinary group equations is classical; for example,
Shmel'kin's nonsingular-system theorem is stated in
\cite[Introduction]{Mikheenko}. Here fixed endomorphisms require
fully invariant vector layers, and the application needs bounds uniform
in their widths. We do not claim a new general theorem on solvability
of group equations or an unrestricted nonlinear-to-linear reduction.

\begin{corollary}[Group codes and exceptional recovery rows]
\notformalized
\label{cor:mc-group-codes}
On the family $E_m$ of \cref{thm:mc-expanding-family}, let $G$ be any
nontrivial finite $2$-group. If all recovery rows have endomorphic
expressions, then
\[
 \log_{|G|}|\mathcal C|\le
 \frac{N_m}{2}+\frac{|E_m|}{2m}
 \le\left(\frac12+\frac{\Delta}{2m}\right)N_m.
\]
More generally, suppose only rows outside a specified $J\subseteq V$
are required to have such expressions. With $S=V\setminus J$ and
$E_m[S]$ the induced incidence relation,
\[
 \log_{|G|}|\mathcal C|\le
 \frac{N_m+|J|}{2}+\frac{|E_m[S]|}{2m}
 \le\frac{N_m+|J|}{2}+\frac{\Delta(N_m-|J|)}{2m}.
\]
These bounds are uniform in group size and expression length.
For $m>\Delta$, a normalized rate at least $1/2+\delta$ requires
\[
 \frac{|J|}{N_m}\ge
       \frac{2\delta-\Delta/m}{1-\Delta/m}.
\]
In particular a sustained excess $\delta>0$ forces
$\liminf_{m\to\infty}|J|/N_m\ge2\delta$.
\end{corollary}
\begin{informalproof}
Restrict the same binary port labels to both ports of coordinates in
$S$. Their matched differences still have rank $m$ and their edge
differences rank one. \Cref{prop:mc-certificate} gives, over $\F_2$
and at every width $t$,
\[
 \dim\ker(I-A)\le t\left(\frac{|S|}{2}
                                      +\frac{|E_m[S]|}{2m}\right).
\]
Apply \cref{thm:mc-group-transfer}. For $J\ne\varnothing$, first
fix the actual coordinates on $J$ and substitute their values as
constants. Each remaining fibre has exactly the required expression
form on $S$. Summing over at most $|G|^{|J|}$ fibres pays for these
coordinates once. Use $|E_m[S]|\le\Delta|S|$ and rearrange to obtain
the final estimate. The empty $S$ case is the trivial full-code bound.
\end{informalproof}

Thus this particular expanding family obstructs more than linear code
spaces: finite-group nonlinearity in the stated language cannot improve
its asymptotic half-rate bound. A construction would need a different
relation or a quantitatively substantial failure of that language.
No small exceptional set or such expression form has been supplied for
arbitrary native cones, and group operations have not been assigned
unit cost as tape operations.
